% ===================================================================
%  圖表第 14 號 — 街。--- 商家。
%
%  Traduction, faite en 2026, en cantonais ecrit d'aujourd'hui, en
%  caracteres traditionnels, sur l'ido de texto/io. Regles de la
%  colonne : voir 10-toubiu-01.tex. Une seule numerotation.
%
%  LA RUE DES COMMERCANTS EST LE TABLEAU DES DEVANTURES, ET LE
%  CANTONAIS Y APPORTE SON VOCABULAIRE D'AFFAIRES ENTIER :
%
%    找換店 (29)  le changeur   Pekin : 兑换处
%    寫字樓 (32)  le bureau     Pekin : 办公室
%    行街 (55)    le voyageur de commerce — celui qui « marche la rue »
%    大律師 (82)  l'avocat      Pekin : 律师
%    舖頭 (15)    la boutique   Pekin : 店铺
%    琴行 (10)    le marchand de musique
%    影樓 (20)    l'atelier du photographe
%    沙井 (68)    la bouche d'egout  Pekin : 窨井
%
%  « 沙井 » est un mot des travaux publics de Hong Kong ; il se lit
%  sur les couvercles de fonte, et le nord ne l'ecrit pas.
%
%  ET LE TABLEAU DONNE ENFIN LE MOT LE PLUS COURT DE LA COLONNE :
%
%    遮 (51)  le parapluie   Pekin : 伞
%
%  Un caractere qui veut dire « couvrir, abriter » et qui, employe
%  seul comme nom, ne designe plus que le parapluie. Le nord garde
%  伞, qui est le dessin de l'objet ; le sud a garde le verbe.
%
%  LES METIERS CONTINUENT EN 佬 — 白鐵佬 (14), 影相佬 (20),
%  打包佬 (19), 掃街佬 (11), 通煙囪佬 (91) — et le vetement prend
%  ses mots : 衣車 la machine a coudre, 車衣 (63) la couture, 西裝
%  (66) le complet, 傢俬 (13) les meubles.
%
%  LA PATISSERIE PARLE ANGLAIS SANS EN AVOIR L'AIR : 忌廉 la creme,
%  批 la tourte, 撻 la tarte. Trois emprunts, trois desserts, et
%  aucun des trois ne se dit au nord.
%
%  ET L'OBSERVATOIRE (1) SE TROUVE ETRE UNE INSTITUTION DE HONG KONG :
%  天文台 est le nom de celle qui donne le temps et hisse les signaux
%  de typhon depuis 1883. Le mot n'a pas ete choisi pour cela, mais
%  il tombe juste.
%
%  QUATRE BLOCS ONT DEMANDE L'INVERSION : la rue et ses bains turcs
%  (t14-02-1), le bureau du cousin (t14-11-1), le jet de la fontaine
%  (t14-14-1), l'imperiale de l'omnibus (t14-15-2).
% ===================================================================

%%K t14-apar-1 apar
\VUsaut{4.0mm}
\VUtitre{15.0pt}{60}{圖表第 14 號}
\VUsaut{3.0mm}

%%K t14-tit-1 sub
\VUpk{11.4pt}{40}{街。--- 商家。}
\VUsaut{3.0mm}

%%K t14-01-1 p
1. --- 我個朋友約翰住喺鄉下，嚟咗我屋企住三日。\VUgras{佢一直}都冇機會見識
大城市，所以我哋成日都用嚟周圍行，我就將佢唔識嘅嘢一樣樣講畀佢聽。

%%K t14-02-1 p
2. --- 我哋第一日去咗參觀個\VUgras{天文台} \textsuperscript{(1)}，佢覺得好有趣。
沿住條\VUgras{街} \textsuperscript{(3)} 返嚟嗰陣——即係有間\VUgras{土耳其浴室}
\textsuperscript{(2)} 嗰條——我哋撞正一場\VUgras{出殯} \textsuperscript{(4)}。
\VUgras{送殯隊伍}最前面行嘅係\VUgras{神職人員，}跟住係\VUgras{靈車，}
\VUgras{棺材}就擺喺入面。好多朋友陪住\VUgras{戴孝}嘅家人行。

%%K t14-02-2 p
隊伍過晒之後，我哋喺間大\VUgras{啤酒館} \textsuperscript{(5)} 門口過對面街；館入面好多
\VUgras{客}喺露台度飲\VUgras{啤酒，}有塊玻璃\VUgras{簷篷} \textsuperscript{(6)} 遮住。

%%K t14-03-1 p
3. --- 有塊\VUgras{盾徽}擺喺間\VUgras{洋酒舖} \textsuperscript{(7)} 同間\VUgras{琴行}
\textsuperscript{(10)} 中間，約翰見到就大為驚訝。佢問我係乜意思；我話嗰度係英國
\VUgras{領事館} \textsuperscript{(8)}，仲指埋個企喺露台嗰位\VUgras{領事} \textsuperscript{(9)} 畀佢睇。

%%K t14-tit-2 sub
\VUpk{10.2pt}{40}{睇舖頭。}
\VUsaut{3.0mm}

%%K t14-04-1 p
4. --- 我哋梗係停低喺啲\VUgras{樂器}櫥窗前面睇——\VUgras{風琴、}\VUgras{管風琴}
同\VUgras{鋼琴；}我哋睇埋啲有插圖嘅樂譜封面，

%%K t14-04-1 p suite
仲有啲鋼琴\VUgras{琴譜；}又欣賞咗個做招牌用嘅鍍金木\VUgras{豎琴。}

%%K t14-05-1 p
5. --- \VUgras{掃街佬} \textsuperscript{(11)} 同架\VUgras{掃街車} \textsuperscript{(12)} 揚起
一陣陣塵，我哋唯有快手行過間\VUgras{傢俬舖} \textsuperscript{(13)}——舖入面啲
\VUgras{傢俬}好靚——又快手行過間\VUgras{白鐵佬} \textsuperscript{(14)} 舖，佢櫥窗擺住
\VUgras{火爐}同\VUgras{燈。}

%%K t14-06-1 p
6. --- 何況行到呢頭，我哋要落返馬路行，因為行人路畀啲包裹塞住——啲包裹啱啱
由架\VUgras{貨車} \textsuperscript{(16)} 度卸落嚟，要搬入間啱租咗嘅\VUgras{舖頭}
\textsuperscript{(15)}。

%%K t14-06-2 p
噉樣我哋就冇得睇\VUgras{造車佬} \textsuperscript{(17)} 啲靚車；不過我哋都停低咗，望住個
\VUgras{打包佬} \textsuperscript{(19)} 幫佢隔籬做緊個木箱——隔籬係間\VUgras{單車汽車行}
\textsuperscript{(18)}。跟住，因為都幾夜嘞，我哋就返屋企。

%%K t14-tit-3 sub
\VUpk{10.2pt}{40}{行街。}
\VUsaut{3.0mm}

%%K t14-07-1 p
7. --- 第二日，我阿媽話不如影張相畀約翰，然之後就叫我陪佢去\VUgras{影相佬}
\textsuperscript{(20)} 度。食完晏，我哋入咗個師傅嘅\VUgras{影樓，}喺嗰度等咗幾耐。
為咗打發時間，我哋就睇下掛喺牆上面啲\VUgras{相} \textsuperscript{(22)}。

%%K t14-07-2 p
輪到我哋嗰陣，做嘢冇幾耐；我個朋友一擺完甫士，喺部\VUgras{影相機}
\textsuperscript{(21)} 前面影完，我哋就落樓。

%%K t14-08-1 p
8. --- 一樓嗰度，我哋隔住間\VUgras{飛髮舖} \textsuperscript{(23)} 開住嘅門，望到佢個
夥計幫緊個客飛髮。

%%K t14-08-2 p
行返出街，我哋經過間\VUgras{煙仔舖} \textsuperscript{(24)}。約翰畀盞紅\VUgras{燈籠}
\textsuperscript{(25)} 同支做\VUgras{招牌} \textsuperscript{(26)} 嘅\VUgras{大煙斗}逗到，行行下就
撞到兩位一路傾偈一路\VUgras{索鼻煙} \textsuperscript{(27)} 嘅老先生。

%%K t14-tit-4 sub
\VUpk{10.2pt}{40}{餅店。}
\VUsaut{3.0mm}

%%K t14-09-1 p
9. --- \VUgras{找換店} \textsuperscript{(29)} 個\VUgras{櫥窗}擺住各國嘅\VUgras{銀紙}同
\VUgras{銀仔，}吸引咗我哋一陣；不過夠鐘食茶點嘞，我哋冇再停，就入咗間\VUgras{餅店}
\textsuperscript{(31)}。

%%K t14-10-1 p
10. --- 舖入面咁多\VUgras{零食：}\VUgras{餅、蜜糖餅、餅乾，}仲有各式各樣嘅
\VUgras{糖，}我哋真係揀到頭都痕。最後約翰揀咗\VUgras{忌廉餅，}我就淨係要咗個細
\VUgras{車厘子撻，}因為食糖會\VUgras{冧我啲牙，}我一啲都唔想成日去見\VUgras{牙醫}
\textsuperscript{(30)}。

%%K t14-tit-5 sub
\VUpk{10.2pt}{40}{打字。}
\VUsaut{3.0mm}

%%K t14-11-1 p
11. --- 我個朋友聽過人講\VUgras{打字機，}但係未見過；為咗畀佢見識下，我哋入咗個
\VUgras{寫字樓} \textsuperscript{(32)}——\VUgras{我表哥} \textsuperscript{(34)} 嗰個。我哋撞到佢
正喺度向個\VUgras{秘書} \textsuperscript{(35)} 講信，個秘書係做\VUgras{速記}嘅。一封信一講完，
個\VUgras{打字員} \textsuperscript{(33)} 就用部機打返一份出嚟。因為唔想阻住我親戚做嘢，
我哋淨係坐咗一陣就走。

%%K t14-12-1 p
12. --- 我表哥個寫字樓樓下，住住一位好大間嘅\VUgras{鐘錶珠寶行}
\textsuperscript{(37)} 老闆，佢仲兼做金器。

%%K t14-12-1 p suite
佢間靚舖入面有好多\VUgras{鐘、掛鐘、錶、}\VUgras{錶鏈，}同貴重嘅\VUgras{珠寶，}
譬如\VUgras{珍珠頸鏈、}金\VUgras{手鈪、耳環}同鑲住\VUgras{寶石}同\VUgras{鑽石}嘅
\VUgras{戒指。}其中一個櫥窗專係擺\VUgras{金器}嘅，仲有一大批銀\VUgras{餐具。}

%%K t14-13-1 p
13. --- 轉個街角嗰陣，我留意到人哋換咗塊\VUgras{門牌} \textsuperscript{(36)}——就擺喺
間屋個\VUgras{門牌號}隔籬。我啱想大大聲講出嚟，就聽到軍樂隊嗰陣開心嘅樂聲。
我哋即刻跑去迎接個軍團，冇諗到\VUgras{佢}就快行過\VUgras{帽舖} \textsuperscript{(39)}
同\VUgras{手套舖} \textsuperscript{(40)} 門口；其實我哋企定喺\VUgras{眼鏡舖}
\textsuperscript{(38)} 個櫥窗前面，一樣睇得清楚——

%%K t14-13-1 p suite
櫥窗入面擺滿\VUgras{夾鼻眼鏡、}\VUgras{單邊眼鏡}同\VUgras{觀劇鏡。}

%%K t14-tit-6 sub
\VUpk{10.2pt}{40}{步操。}
\VUsaut{3.0mm}

%%K t14-14-1 p
14. --- \VUgras{軍團} \textsuperscript{(41)} 最前面行嘅係\VUgras{鼓手長}
\textsuperscript{(42)}，後面跟住啲\VUgras{鼓手} \textsuperscript{(43)}、啲\VUgras{號手}
\textsuperscript{(44)}，同成隊\VUgras{樂隊。}\VUgras{將軍} \textsuperscript{(45)} 同佢副官
喺廣場度，停咗喺道\VUgras{水柱} \textsuperscript{(49)} 隔籬——即係個\VUgras{噴泉}
\textsuperscript{(48)} 嗰道——睇\VUgras{步操。}

%%K t14-14-2 p
\VUgras{啲參謀} \textsuperscript{(46)}、\VUgras{上校}同\VUgras{中校}舉劍向佢敬禮。
\VUgras{啲少校}企喺自己啲\VUgras{營} \textsuperscript{(47)} 前面，每個\VUgras{上尉}指揮住
自己個\VUgras{連；}啲\VUgras{中尉、少尉}同\VUgras{士官}就企返自己慣常嘅位，
喺自己啲兵隔籬。

%%K t14-15-1 p
15. --- 好多路人喺個\VUgras{十字路口} \textsuperscript{(50)} 度圍埋一堆；啲商家
——賣\VUgras{遮嘅} \textsuperscript{(51)}、\VUgras{打錫佬} \textsuperscript{(53)}、\VUgras{槍械舖}
\textsuperscript{(54)}——都行咗出嚟睇\VUgras{啲兵。}所有車，\VUgras{出租馬車}
\textsuperscript{(52)}、\VUgras{搬運車} \textsuperscript{(57)}，連\VUgras{灑水車} \textsuperscript{(56)}
都埋咗埋行人路邊。

%%K t14-tit-7 sub
\VUpk{10.2pt}{40}{街頭一幕。}
\VUsaut{3.0mm}

%%K t14-15-2 p
一個\VUgras{行街} \textsuperscript{(55)} 手上攞住佢啲\VUgras{辦板箱，}快步過咗對面街；
坐喺\VUgras{上層} \textsuperscript{(59)} 嘅人——\VUgras{公共馬車} \textsuperscript{(58)} 個上層
——都擰轉身睇熱鬧。有個可憐嘅\VUgras{街頭歌手} \textsuperscript{(60)} 帶住佢部
\VUgras{手搖風琴} \textsuperscript{(61)}，唯有停低唔唱，因為啲鼓聲蓋過晒佢把\VUgras{歌聲。}

%%K t14-16-1 p
16. --- 個軍團行到間\VUgras{布疋舖} \textsuperscript{(64)} 前面嗰陣，啲\VUgras{車衣女}
\textsuperscript{(63)} 同啲\VUgras{裁縫} \textsuperscript{(62)} 都放低部\VUgras{衣車}同手上啲嘢，
走去窗口望。\VUgras{個老闆} \textsuperscript{(65)} 啱啱先將幾套新\VUgras{西裝}
\textsuperscript{(66)} 掛上\VUgras{櫥窗，}都要將擺喺行人路上面嗰啲\VUgras{呢絨}
\textsuperscript{(67)} 同\VUgras{布疋}搬返入舖——布就攤喺個\VUgras{沙井}
\textsuperscript{(68)} 隔籬，佢驚啲人逼親打瀉。連個老\VUgras{走鬼檔}
\textsuperscript{(69)} 都要行入條走廊先做得完單生意——有個看更婆正同佢
講緊價買襪。

%%K t14-16-2 p
我哋跟住個軍團一直行到兵營；\VUgras{我哋}呢日過得好開心，就趕返屋企食晚飯。

%%K t14-tit-8 sub
\VUpk{10.2pt}{40}{各行各業。}
\VUsaut{3.0mm}

%%K t14-17-1 p
17. --- 第二日，我叔叔話不如帶我哋去參觀嗰份報紙嘅印刷廠。我哋一口應承。

%%K t14-17-2 p
個\VUgras{印刷佬}同時仲係\VUgras{書商} \textsuperscript{(70)}\VUgras{（即係賣書嘅）、}
\VUgras{出版商、}\VUgras{文具商}同\VUgras{釘裝師傅。}佢喺一樓開咗份報紙嘅
\VUgras{編輯部} \textsuperscript{(71)}，仲有個\VUgras{製版房，}入面做嘢嘅係啲
\VUgras{石印師傅} \textsuperscript{(72)} 同啲\VUgras{銅版師傅} \textsuperscript{(73)}；最後係
\VUgras{排字房，}入面係啲\VUgras{排字工} \textsuperscript{(74)}。真真正正嘅\VUgras{印刷房}
\textsuperscript{(75)}、啲\VUgras{印刷機}同各式機器，就喺地下。

%%K t14-tit-9 sub
\VUpk{10.2pt}{40}{約翰問長問短。}
\VUsaut{3.0mm}

%%K t14-18-1 p
18. --- 行出印刷廠嗰陣，約翰好驚訝噉停低喺一個裝喺柱上面嘅細玻璃盒前面
——就喺間\VUgras{鈕扣舖} \textsuperscript{(77)} 對面——問嗰樣嘢做乜用。我叔叔同佢解釋，
話嗰個係\VUgras{火警鐘} \textsuperscript{(76)}。其實城入面乜都令我個朋友大開眼界：由
最簡單嘅\VUgras{石水池} \textsuperscript{(78)}、輕便嘅\VUgras{三輪送貨車}
\textsuperscript{(80)}——踩車嗰個係着住制服嘅\VUgras{後生} \textsuperscript{(79)}
——一直去到架\VUgras{郵車} \textsuperscript{(81)}。

%%K t14-19-1 p
19. --- 我叔叔行開咗一陣去同個\VUgras{稅務員} \textsuperscript{(83)} 講嘢——嗰位停低咗
同一位\VUgras{大律師} \textsuperscript{(82)} 同一位\VUgras{律師} \textsuperscript{(84)} 傾偈——我哋
就自己用眼跟住街上面啲來來往往。各式各樣嘅人喺我哋面前行過：一個
\VUgras{餅店夥計} \textsuperscript{(85)} 用個籮載住個靚\VUgras{批，}跟喺一個\VUgras{賣衫佬}
\textsuperscript{(86)} 後面；一個\VUgras{點燈佬} \textsuperscript{(87)} 同一個\VUgras{煤氣工}
\textsuperscript{(88)} 傾緊偈；一位\VUgras{修女} \textsuperscript{(89)} 拖住個孤兒女仔，正返緊
自己間修院。

%%K t14-tit-10 sub
\VUpk{10.2pt}{40}{返屋企路上。}
\VUsaut{3.0mm}

%%K t14-20-1 p
20. --- 我叔叔行返埋嚟嗰一刻，約翰笑到彎晒腰：佢見到兩個\VUgras{通煙囪佬}
\textsuperscript{(91)} 成身煤煙，塊面又污糟，件衫又古怪。

%%K t14-21-1 p
21. --- 我想喺個\VUgras{花檔女} \textsuperscript{(92)} 度買盆花畀我阿媽，不過我叔叔提我，
話盆嘢太重唔好攞，不如買啲\VUgras{橙}好啲。我哋行埋去個\VUgras{女檔主}
\textsuperscript{(94)} 度；等個\VUgras{家庭教師} \textsuperscript{(93)}——佢正喺度幫自己個學生揀
——買完，我哋就叫佢招呼我哋。

%%K t14-22-1 p
22. --- 返屋企路上，我哋撞到幾個熟人：我哋家一位老朋友，佢手挽住自己個
\VUgras{女伴} \textsuperscript{(95)}；管橋樑同道路嘅\VUgras{工程師} \textsuperscript{(96)}；仲有
我一個表姐同佢個\VUgras{保姆} \textsuperscript{(98)}。

%%K t14-22-2 p
跟住我哋又撞到我阿媽個\VUgras{帽舖女老闆} \textsuperscript{(97)}，同兩位唔係本城人嘅
\VUgras{修士} \textsuperscript{(99)}——佢哋行埋嚟問我叔叔去火車站點行。

\VUsaut{6.0mm}
\VUfilet{22mm}
